%%%PAGE 148 rot=0 legibility=good summary=电路化简：(1)桥式电路R3电流为零无法化简为串并联；(2)惠斯通电桥及三角形网络求R_AB；(3)节点法与等电势原理，及正方形带对角线电阻网络求R_AB
% 页面背面有淡淡的透印字迹(忽略)
%%%BLOCK id=148-01 ch=10 sec=电路简化·惠斯通电桥 kind=knowledge alt=none cont=none y=0.0-0.16
\textbf{电路化简}

\noindent (1)
%FIG p148_a 0.08 0.013 0.385 0.155
\notefig[0.3]{figures/p148_a.png}
\noindent 该图无法化简为串并联
%%%END
%%%BLOCK id=148-02 ch=10 sec=电路简化·惠斯通电桥 kind=knowledge alt=none cont=none y=0.17-0.31
\noindent (2)　惠斯通电桥　$\dfrac{R_1}{R_2}=\dfrac{R_4}{R_5}$　对角线乘积相等
%FIG p148_b 0.065 0.215 0.41 0.31
\notefig[0.35]{figures/p148_b.png}
%%%END
%%%BLOCK id=148-03 ch=10 sec=电路简化·惠斯通电桥 kind=problem alt=none cont=none y=0.31-0.49
% 四幅图之间有一条手绘长曲线相连；图内文字(求R_AB、各电阻编号、A/B等)已在图中，右下角的 R_AB 算式在图内也有，下方单独转录
\begin{problem}
求 $R_{AB}$
%FIG p148_c 0.08 0.32 0.98 0.49
\notefig[0.85]{figures/p148_c.png}
\begin{solution}
$R_{AB}=\dfrac{4\times2}{4+2}+1=\dfrac{7}{3}\,\Omega$
\end{solution}
\end{problem}
%%%END
%%%BLOCK id=148-04 ch=10 sec=电路简化·节点法 kind=knowledge alt=none cont=none y=0.50-0.80
% 原稿用手绘括线表示层级：“节点法”→“由2根及以上导线汇交而成的点”；“等电势原理”→“在一根导线上所有点电势相等”；“等电势原理”中的“电”字叠写在“等”“势”之间，“点”后有两个小点
\noindent \textbf{(3)　节点法}
\begin{itemize}
\item 由2根及以上导线汇交而成的点..
\item 等电势原理：电压$=$电势降
\begin{itemize}
\item 在一根导线上所有点电势相等
\end{itemize}
\end{itemize}
%FIG p148_d 0.06 0.70 0.455 0.79
\notefig[0.4]{figures/p148_d.png}
%FIG p148_e 0.485 0.70 0.73 0.805
\notefig[0.25]{figures/p148_e.png}
%%%END
%%%BLOCK id=148-05 ch=10 sec=电路简化·节点法 kind=problem alt=none cont=none y=0.79-1.00
% 左下方电路图各角点标有字母 A/a、B/b、M/b 等，已在图内
\begin{problem}
%FIG p148_f 0.02 0.79 0.33 0.977
\notefig[0.3]{figures/p148_f.png}
求 $R_{AB}$
\begin{solution}
%FIG p148_g 0.42 0.83 0.715 0.97
\notefig[0.3]{figures/p148_g.png}
\[ \dfrac{\frac{3}{2}\,\frac{1}{2}}{\frac{3}{2}+\frac{1}{2}}=\dfrac{3}{8}R \]% CHECK: 分子两分数间原稿无运算符，按上下文应为乘积
\end{solution}
\end{problem}
%%%END
%%%PAGEEND 148
